\centerline{\bf \Huge Квадратные трёхчлены}

\begin{flushright}
        \scriptsize{Эпизод 10. Ныряльщик в магме }
\end{flushright}

\problem{1}
Найдите многочлен второй степени, если известно, что его корни равны $-\frac{4}{7}$ и $\frac{5}{3}$, а свободный член равен -2 .

\problem{2}
Найдите все пары квадратных трёхчленов $x^2+a x+b$, $x^2+c x+d$ такие, что $a$ и $b$ - корни второго трёхчлена, $c$ и $d$ - корни первого.

\problem{3}
Даны квадратные трёхчлены $x^2+a x+b$, $x^2+c x+d$ и $x^2+e x+f$. Оказалось, что любые два из них имеют общий корень, но все три общего корня не имеют. Докажите, что выполнены ровно два неравенства из следующих трёх:

$$
\begin{aligned}
& \frac{a^2+c^2-e^2}{4}>b+d-f \\
& \frac{c^2+e^2-a^2}{4}>d+f-b \\
& \frac{e^2+a^2-c^2}{4}>f+b-d
\end{aligned}
$$

\problem{4}
Существуют ли такие попарно различные числа $a, b$ и $c$, что число $a$ является корнем квадратного трёхчлена $x^2-2 b x+c^2$, число $b$ является корнем квадратного трёхчлена $x^2-2 c x+a^2$, а число $c$ является корнем квадратного трёхчлена $x^2-2 a x+b^2$ ?

\problem{5}
Числа $a$ и $b$ таковы, что каждый из двух квадратных трёхчленов $x^2+a x+b$ и $x^2+b x+a$ имеет по два различных корня, а произведение этих трёхчленов имеет ровно три различных корня. Найдите все возможные значения суммы этих трёх корней.

\problem{6}
Различные числа $a, b$ и $c$ таковы, что уравнения $x^2+a x+1=0$ и $x^2+b x+c=0$ имеют общий действительный корень. Кроме того, общий действительный корень имеют уравнения $x^2+x+a=0$ и $x^2+c x+b=0$. Найдите сумму $a+b+c$.